\chapter{Business Model}
\label{ch:business}

The challenge requires a business-model note, and the honest version of one for the final
report is not a revenue projection. It is an answer to three questions: who has this
problem badly enough to change behaviour, what exactly they receive, and how the thing pays
for itself and survives its provider.

\section{Target users}

\begin{table}[htbp]
\centering
\caption{Three segments, and the reason each one is reachable without a sales team.}
\label{tab:segments}
\footnotesize
\begin{tabularx}{\textwidth}{@{}p{2.9cm}X p{4.2cm}@{}}
\toprule
\textbf{Segment} & \textbf{Why they feel the problem} & \textbf{Entry point} \\
\midrule
University tutoring and revision circles & Small groups already meet; the failure is
  coordination, and nobody is paid to fix it & IEEE student branches and campus clubs \\
Training centres and bootcamps & Mixed-level cohorts where one learner dominates and the
  trainer cannot watch every group & The facilitator board is the trainer's tool \\
Self-learning communities & Motivated learners with no facilitator at all; they need
  structure more than content & Public circles, join by code or QR \\
\bottomrule
\end{tabularx}
\end{table}

The three segments share a property that makes adoption plausible: the product enters an
existing ritual rather than asking for a new one. A revision circle already meets; HALQA
changes how the meeting runs, not whether it happens.

\section{Value proposition}

The proposition is stated in terms its buyer can check.

\begin{itemize}[leftmargin=1.4em,itemsep=2pt]
  \item \textbf{Fair turns without extra facilitation.} The window and the \emph{Share}
        ring make participation visible, so a group can correct itself without an adult in
        the room.
  \item \textbf{Fewer wrong answers spreading.} The Auditor and poll-first behaviour
        intervene before a misconception becomes consensus, which is the expensive failure
        in peer learning.
  \item \textbf{Memory of the session.} The revision capsule and the quiz bank turn a
        conversation into something revisable, at no extra effort.
  \item \textbf{Bounded cost.} Token spend is a documented policy with a measured ceiling,
        so the cost of a session is predictable rather than proportional to enthusiasm.
  \item \textbf{Visibility for the person responsible.} The board shows which circle needs a
        human, which is the only feature a training centre will pay for.
\end{itemize}

\section{Adoption path}

\paragraph{Step 1 --- bottom-up pilot.} IEEE student branches and campus clubs already in
the WIE ACT network run HALQA for one revision cycle. The two things learned from a pilot
are not usage statistics but \emph{pack demand}: which course the cohort actually revises.
That feeds the seed knowledge base, which is the real asset.

\paragraph{Step 2 --- department licensing.} Tutoring programmes buy per-cohort access, and
the board view is the stakeholder-facing artefact. A department does not buy a chatbot; it
buys evidence that its tutoring hours are working.

\paragraph{Step 3 --- training-centre integration.} Capsule and quiz exports are designed to
drop into an existing LMS workflow rather than to replace it. The export is the integration
surface, which avoids an enterprise sales cycle beyond this submission.

\section{Sustainability beyond the challenge}

Two engineering decisions from earlier chapters are also business decisions, and this is
where they pay.

\begin{table}[htbp]
\centering
\caption{Engineering decisions and their commercial consequence.}
\label{tab:sustain}
\footnotesize
\begin{tabularx}{\textwidth}{@{}p{4.0cm}X@{}}
\toprule
\textbf{Decision} & \textbf{Commercial consequence} \\
\midrule
One provider adapter & The product survives a provider switch or a migration to
  self-hosted open-weight models without a feature freeze \\
Two-tier routing with a measured ceiling & Cost per session is bounded and predictable;
  the larger tier is a deliberate purchase, not a default \\
TF-IDF retrieval, no embedding service & No per-query inference cost on the grounded path,
  and retrieval keeps working when the provider does not \\
Zero-GPU workload & Running cost is close to hosting cost; the margin does not depend on
  negotiated GPU capacity \\
Seed-format knowledge bases & Each new course is a JSON pack plus a manifest row, so
  content can be contributed by adopters instead of authored centrally \\
\bottomrule
\end{tabularx}
\end{table}

\section{Cost model}

The cost structure is unusually simple because the expensive variable is a call count rather
than a token count.

\begin{table}[htbp]
\centering
\caption{Cost drivers and the control that bounds each.}
\label{tab:cost}
\footnotesize
\begin{tabularx}{\textwidth}{@{}p{3.3cm}p{3.4cm}X@{}}
\toprule
\textbf{Driver} & \textbf{Bound} & \textbf{Control} \\
\midrule
Compact calls & 20 per session & Governor ledger, atomic reserve \\
Larger calls & 5 per session & Separate cap; refusal downgrades to C \\
Class-level spend & 400 calls per circle per day & Daily ceiling, the binding limit \\
Repeat questions & Served from cache & Semantic cache with a metered hit rate \\
Hosting & One API process, one Mongo & No GPU, no vector store, no queue service \\
\bottomrule
\end{tabularx}
\end{table}

\section{Roadmap}

\begin{table}[htbp]
\centering
\caption{Roadmap by horizon. Items are ordered by what unblocks the next thing.}
\label{tab:roadmap}
\footnotesize
\begin{tabularx}{\textwidth}{@{}p{2.4cm}X@{}}
\toprule
\textbf{Horizon} & \textbf{Work} \\
\midrule
Rehearsal and pilot & Real-user testing in two contexts (university and self-learning),
  with the four-language interaction evidence recorded alongside the automated evals \\
Horizon 1 & Teacher dashboard exports, additional course packs, tuned bandit parameters from
  observed learners, consent capture \\
Horizon 2 & Voice interaction; collaborative quiz battles at scale; a second provider
  adapter as a resilience measure rather than a feature \\
\bottomrule
\end{tabularx}
\end{table}

\begin{note}[What this chapter does not claim]
There is no revenue projection, no pricing table and no market-size figure. For this
submission those numbers would be invented, and an invented number in a business section
costs more credibility than an absent one.
The horizons below are delivery sequence, not financial forecast.
\end{note}
